
%\needspace{4\baselineskip}
\color{uniblue}{\subsection{Логика связи (ЛС)}}
\color{black}

\begin{enumerate}[label=\arabic{section}.\arabic{subsection}.\arabic{enumi}, labelsep=4pt, leftmargin=0pt, itemindent=57pt, itemsep=0pt, parsep=5pt] % если что можно крутить itemsep

%%%%%%%%%%%%%%%%%%%%%%%%%%%%%%%%%%%%%%%%%%%%%%%%%%%%%%%%%%%%%%%%%%%%%%  ОБЩИЕ СВЕДЕНИЯ %%%%%%%%%%%%%%%%%%%%%%%%%%%%%%%%%%%%%%%%%%%%%%%%%%
\item Общие сведения
\begin{enumerate}[label=\arabic{section}.\arabic{subsection}.\arabic{enumi}.\arabic*, labelsep=4pt, leftmargin=0em, itemindent=65pt, parsep=0pt]

\item
В устройстве предусмотрена возможность организации приема и передачи сигналов и команд телеотключения и телеускорения для обеспечения быстрого селективного отключения КЗ на защищаемой ЛЭП, которые также могут называться сигналами высокочастотного телеотключения (ВЧТО).

\item
Логика связи может быть оперативно выведена из работы от внешнего сигнала \textit{«Вывод ЛС»}, назначенного на соответствующий дискретный вход или функциональную клавишу. Из действия выводятся как прием, так и передача всех команд телеотключения и телеускорения, в том числе эхо-логика и ЛОСК.

\end{enumerate}

%%%%%%%%%%%%%%%%%%%%%%%%%%%%%%%%%%%%%%%%%%%%%%%%%%%%%%%%%%%%%%%%%%%%%%  ТЕЛЕОТКЛЮЧЕНИЕ %%%%%%%%%%%%%%%%%%%%%%%%%%%%%%%%%%%%%%%%%%%%%%%%%%
\item Телеотключение
\begin{enumerate}[label=\arabic{section}.\arabic{subsection}.\arabic{enumi}.\arabic*, labelsep=4pt, leftmargin=0em, itemindent=65pt, parsep=0pt]

\item
Формирование сигнала на отключение выключателя при приеме сигнала телеотключения (\textit{«Прием Тоткл»}) может осуществляться с подтверждением от следующих сигналов:

\begin{itemize}
\item \textit{«БК-I-м»} (контроль вводится программным переключателем \textbf{«Контр. Тоткл от БК-I-м»});
\item отключенное состояние ЛЭП (контроль вводится программным переключателем \textbf{«Контр. Тоткл от РПО»});
\item пуск одной из ступеней ДЗ-ФФ или ДЗ-ФЗ (ступени выбираются программными переключателями \textbf{«Контр. Тоткл от ДЗ-ФФ»} и \textbf{«Контр. Тоткл от ДЗ-ФЗ»});
\item срабатывание ИО ТНЗНП выбранной ступени (выбор осуществляется программным переключателем \textbf{«Контр. Тоткл от ТНЗНП»});
\end{itemize}

\item
Функциональная схема логики приема/выдачи сигнала отключения приведена на рисунке \ref{pict_lo:ls1}.

\begin{figure}[H]
\centering
\includegraphics[width=0.9\textwidth,height=0.9\textheight,keepaspectratio]{\fbpath/561.1711 Логика связи/img/logic_Totkl.pdf}
\captionof{figure}{Функциональная схема логики телеотключения}\label{pict_lo:ls1}
\end{figure}

\item 
Если программный ключ \textbf{«Прием Тоткл с контр.»} выведен, то сигнал на отключение выключателя при приеме сигнала телеотключения формируется без контролей.

\item
Выдача сигнала телеотключения \textit{«Выход Тоткл»} на противоположный конец линии осуществляется при срабатывании ЗНР, срабатывании внутреннего УРОВ или по сигналу срабатывания внешнего УРОВ.

\item
Уставки телеотключения представлены в таблице \ref{desc_lo:tbl1}.

\small
\begin{longtable}{|>{\centering\arraybackslash}m{5.3cm}|>{\centering\arraybackslash}m{3.3cm}|>{\centering\arraybackslash}m{4.2cm}|>{\centering\arraybackslash}m{1.8cm}|>{\centering\arraybackslash}m{1cm}|}
\caption{Параметры для настройки телеотключения\hfill\vspace{-0.5\baselineskip}}\label{desc_lo:tbl1}\\ 
\hline
\rowcolor{gray!30}
Параметр & Условное обозначение на схеме & Значение/ Диапазон & Единица измерения & Шаг \\ 
\hline
\endfirsthead
\caption*{\hspace{3pt}\emph{Продолжение таблицы \ref{desc_lo:tbl1}\hfill\vspace{-0.5\baselineskip}}} \\ % сделано по ГОСТ 2.105 п.6.8.7
\hline
\rowcolor{gray!30}
Параметр & Условное обозначение на схеме & Значение/ Диапазон & Единица измерения & Шаг \\ 
%\hline
\endhead

%===m>Логика связи|Телеотключение
\raggedright {'Description': ''} & \centering Programmnyjklyuch2s & \centering Введено \\ Выведено & \centering --- & \centering \arraybackslash --- \\
\hline
\raggedright {'Description': ''} & \centering Programmnyjklyuch1s & \centering Введено \\ Выведено & \centering --- & \centering \arraybackslash --- \\
\hline
\raggedright {'Description': ''} & \centering PereklyuchatelVkh2boolconst1S & \centering Выведено \\ Введено & \centering --- & \centering \arraybackslash --- \\
\hline
\raggedright {'Description': ''} & \centering PereklyuchatelVkh7boolconst1S & \centering Откл. \\ 1 ступ. \\ 2 ступ. \\ 3 ступ. \\ 4 ступ. \\ 5 ступ. \\ 6 ступ. & \centering --- & \centering \arraybackslash --- \\
\hline
\raggedright {'Description': ''} & \centering PereklyuchatelVkh4boolconst1S & \centering Откл. \\ 1 ступ. \\ 2 ступ. \\ 3 ступ. & \centering --- & \centering \arraybackslash --- \\
\hline
\raggedright {'Description': ''} & \centering PereklyuchatelVkh7boolconst2S & \centering Откл. \\ 1 ступ. \\ 2 ступ. \\ 3 ступ. \\ 4 ступ. \\ 5 ступ. \\ 6 ступ. & \centering --- & \centering \arraybackslash --- \\
\hline
\raggedright {'Description': ''} & \centering PereklyuchatelVkh2boolconst2S & \centering Выведено \\ Введено & \centering --- & \centering \arraybackslash --- \\
\hline
%===m
\end{longtable}
\normalsize

\end{enumerate}

%%%%%%%%%%%%%%%%%%%%%%%%%%%%%%%%%%%%%%%%%%%%%%%%%%%%%%%%%%%%%%%%%%%%%%  ТЕЛЕУСКОРЕНИЕ ТРЕХ ФАЗ %%%%%%%%%%%%%%%%%%%%%%%%%%%%%%%%%%%%%%%%%%%%%%%%%%
\item Телеускорение отключения трех фаз
\begin{enumerate}[label=\arabic{section}.\arabic{subsection}.\arabic{enumi}.\arabic*, labelsep=4pt, leftmargin=0em, itemindent=65pt, parsep=0pt]

\item
Прием сигнала телеускорения отключения трёх фаз (ТУ ОТФ) может осуществляться с контролем срабатывания:

\begin{itemize}
    \item направленного РС-ФФ-2 с охватом начала координата (вводится программным ключом \textbf{«Контр. РС-ФФ-2 с охв.»});
    \item направленного РС-ФФ-3 (вводится программным ключом \textbf{«Контр. РС-ФФ-3»});
    \item направленного РС-ФФ-4 (вводится программным ключом \textbf{«Контр. РС-ФФ-4»});
    \item направленного РС-ФЗ-2 (вводится программным ключом \textbf{«Контр. РС-ФЗ-2»});
    \item ИО четвертой ступени ТНЗНП (вводится программным ключом \textbf{«Контр. ИО ТНЗНП-4»}).
\end{itemize}

\item
В зависимости от положения программного переключателя \textbf{«БК ТУ ОТФ»} выбранные ступени РС для подтверждения ТУ ОТФ могут контролироваться от БК-I-б, БК-I-м или БК-Z. Для ИО ТНЗНП-4 программным ключом \textbf{«Контр. РНМНП-Р»} может быть введен контроль направленности.

\item
Выдержка времени на срабатывание ТУ ОТФ регулируется уставкой \textbf{«Тср ТУ ОТФ»}.

\item
Функциональная схема логики ТУ ОТФ приведена на рисунке \ref{pict_lo:ls2}.

\item
Выдача разрешающего сигнала ТУ ОТФ \textit{«Выход ТУ ОТФ»} на противоположный конец линии осуществляется при срабатывании внутренних защит устройства, кроме ТЗО.

\begin{figure}[H]
\centering
\includegraphics[width=0.9\textwidth,height=0.9\textheight,keepaspectratio]{\fbpath/561.1711 Логика связи/img/logic_TU_OTF.pdf}
\captionof{figure}{Функциональная схема логики ТУ ОТФ}\label{pict_lo:ls2}
\end{figure}

\item
Уставки ТУ ОТФ представлены в таблице \ref{desc_lo:tbl2}.

\small
\begin{longtable}{|>{\centering\arraybackslash}m{5.3cm}|>{\centering\arraybackslash}m{3.3cm}|>{\centering\arraybackslash}m{4.2cm}|>{\centering\arraybackslash}m{1.8cm}|>{\centering\arraybackslash}m{1cm}|}
\caption{Параметры для настройки ТУ ОТФ\hfill\vspace{-0.5\baselineskip}}\label{desc_lo:tbl2}\\ 
\hline
\rowcolor{gray!30}
Параметр & Условное обозначение на схеме & Значение/ Диапазон & Единица измерения & Шаг \\ 
\hline
\endfirsthead
\caption*{\hspace{3pt}\emph{Продолжение таблицы \ref{desc_lo:tbl2}\hfill\vspace{-0.5\baselineskip}}} \\ % сделано по ГОСТ 2.105 п.6.8.7
\hline
\rowcolor{gray!30}
Параметр & Условное обозначение на схеме & Значение/ Диапазон & Единица измерения & Шаг \\ 
%\hline
\endhead

%===m>Логика связи|ТУ ОТФ
\raggedright {'Description': ''} & \centering NezavisimoeRVnasrabatyvanie1Tvremya & \centering 0 ... 100000 & \centering  & \centering \arraybackslash 1 \\
\hline
\raggedright {'Description': ''} & \centering Programmnyjklyuch1s & \centering Введено \\ Выведено & \centering --- & \centering \arraybackslash --- \\
\hline
\raggedright {'Description': ''} & \centering Programmnyjklyuch2s & \centering Введено \\ Выведено & \centering --- & \centering \arraybackslash --- \\
\hline
\raggedright {'Description': ''} & \centering Programmnyjklyuch3s & \centering Введено \\ Выведено & \centering --- & \centering \arraybackslash --- \\
\hline
\raggedright {'Description': ''} & \centering Programmnyjklyuch4s & \centering Введено \\ Выведено & \centering --- & \centering \arraybackslash --- \\
\hline
\raggedright {'Description': ''} & \centering PereklyuchatelVkh4boolconst1S & \centering Откл. \\ БК-I-б \\ БК-I-м \\ БК-Z & \centering --- & \centering \arraybackslash --- \\
\hline
\raggedright {'Description': ''} & \centering Programmnyjklyuch5s & \centering Введено \\ Выведено & \centering --- & \centering \arraybackslash --- \\
\hline
\raggedright {'Description': ''} & \centering PereklyuchatelVkh2boolconst1S & \centering Выведено \\ Введено & \centering --- & \centering \arraybackslash --- \\
\hline
\raggedright {'Description': ''} & \centering PereklyuchatelVkh2boolconst2S & \centering Выведено \\ Введено & \centering --- & \centering \arraybackslash --- \\
\hline
%===m
\end{longtable}
\normalsize

\end{enumerate}

%%%%%%%%%%%%%%%%%%%%%%%%%%%%%%%%%%%%%%%%%%%%%%%%%%%%%%%%%%%%%%%%%%%%%%  ТЕЛЕУСКОРЕНИЕ ДЗ %%%%%%%%%%%%%%%%%%%%%%%%%%%%%%%%%%%%%%%%%%%%%%%%%%
\item Телеускорение ступеней ДЗ
\begin{enumerate}[label=\arabic{section}.\arabic{subsection}.\arabic{enumi}.\arabic*, labelsep=4pt, leftmargin=0em, itemindent=65pt, parsep=0pt]

\item
Для реализации логики работы ДЗ с использованием разрешающих сигналов предусмотрены прием и выдача сигнала телеускорения соответствующей ступени ДЗ (ТУ ДЗ).

\item
Прием сигнала ТУ ДЗ осуществляется с контролем пуска выбранных ступеней ДЗ-ФФ и ДЗ-ФЗ, направленных в линию -- выбор осуществляется программными переключателями \textbf{«Прямонапр. ступень ДЗ-ФФ»} и \textbf{«Прямонапр. ступень ДЗ-ФЗ»}. Прямая направленность задается в соответствующих настройках самой ступени.

\item
Для исключения неселективной работы ТУ ДЗ применяется блокировка при реверсе мощности. Для выявления режима реверса мощности используются обратнонаправленные ступени ДЗ-ФФ и ДЗ-ФЗ -- выбор ступеней осуществляется программными переключателями \textbf{«Обратнонапр. ступень ДЗ-ФФ»} и \textbf{«Обратнонапр. ступень ДЗ-ФЗ»}. Обратная направленность задается в соответствующих настройках самой ступени.

\item
Выдержка времени \textbf{«Тфикс об. ступени»} обеспечивает фиксацию устойчивого пуска обратнонаправленной ступени ДЗ при КЗ на параллельной линии. Его значение должно быть меньше полного времени отключения повреждения на параллельной линии. Выдержка времени \textbf{«Тпрод об. ступени»} обеспечивает продление сигнала о пуске обратнонаправленной ступени до срабатывания прямонаправленной ступени. По одновременному наличию продленного сигнала пуска обратнонаправленной ступени и появлению пуска прямонаправленной ступени фиксируется режим реверса мощности. 

\item
Минимальное время действия блокировки при реверсе мощности определяется выдержкой времени \textbf{«Тв блок. ТУ ДЗ»}.

\item
Выдержка времени на срабатывание ТУ ДЗ регулируется уставкой \textbf{«Тср ТУ ДЗ»}.

\item
Разрешающий сигнал \textit{«Выход ТУ ДЗ»} на противоположный конец линии выдается с контролем пуска прямонаправленной ступени ДЗ и с контролем реверса мощности.

\item
Функциональная схема логики ТУ ДЗ приведена на рисунке \ref{pict_lo:tu_dz}.

\begin{figure}[H]
\centering
\includegraphics[width=0.9\textwidth,height=0.9\textheight,keepaspectratio]{\fbpath/561.1711 Логика связи/img/logic_TU_DZ.pdf}
\captionof{figure}{Функциональная схема логики ТУ ДЗ}\label{pict_lo:tu_dz}
\end{figure}

\item
Уставки ТУ ДЗ представлены в таблице \ref{desc_lo:tbl3}.

\small
\begin{longtable}{|>{\centering\arraybackslash}m{5.3cm}|>{\centering\arraybackslash}m{3.3cm}|>{\centering\arraybackslash}m{4.2cm}|>{\centering\arraybackslash}m{1.8cm}|>{\centering\arraybackslash}m{1cm}|}
\caption{Параметры для настройки ТУ ДЗ\hfill\vspace{-0.5\baselineskip}}\label{desc_lo:tbl3}\\ 
\hline
\rowcolor{gray!30}
Параметр & Условное обозначение на схеме & Значение/ Диапазон & Единица измерения & Шаг \\ 
\hline
\endfirsthead
\caption*{\hspace{3pt}\emph{Продолжение таблицы \ref{desc_lo:tbl3}\hfill\vspace{-0.5\baselineskip}}} \\ % сделано по ГОСТ 2.105 п.6.8.7
\hline
\rowcolor{gray!30}
Параметр & Условное обозначение на схеме & Значение/ Диапазон & Единица измерения & Шаг \\ 
%\hline
\endhead

%===m>Логика связи|ТУ ДЗ
\raggedright {'Description': ''} & \centering NezavisimoeRVnasrabatyvanie5Tvremya & \centering 0 ... 100000 & \centering  & \centering \arraybackslash 1 \\
\hline
\raggedright {'Description': ''} & \centering PereklyuchatelVkh7boolconst1S & \centering Откл. \\ 1 ступ. \\ 2 ступ. \\ 3 ступ. \\ 4 ступ. \\ 5 ступ. \\ 6 ступ. & \centering --- & \centering \arraybackslash --- \\
\hline
\raggedright {'Description': ''} & \centering NezavisimoeRVnavozvrat1Tvremya & \centering 0 ... 100000 & \centering  & \centering \arraybackslash 1 \\
\hline
\raggedright {'Description': ''} & \centering NezavisimoeRVnasrabatyvanie1Tvremya & \centering 0 ... 100000 & \centering  & \centering \arraybackslash 1 \\
\hline
\raggedright {'Description': ''} & \centering NezavisimoeRVnasrabatyvanie2Tvremya & \centering 0 ... 100000 & \centering  & \centering \arraybackslash 1 \\
\hline
\raggedright {'Description': ''} & \centering PereklyuchatelVkh7boolconst2S & \centering Откл. \\ 1 ступ. \\ 2 ступ. \\ 3 ступ. \\ 4 ступ. \\ 5 ступ. \\ 6 ступ. & \centering --- & \centering \arraybackslash --- \\
\hline
\raggedright {'Description': ''} & \centering PereklyuchatelVkh4boolconst1S & \centering Откл. \\ 1 ступ. \\ 2 ступ. \\ 3 ступ. & \centering --- & \centering \arraybackslash --- \\
\hline
\raggedright {'Description': ''} & \centering PereklyuchatelVkh4boolconst2S & \centering Откл. \\ 1 ступ. \\ 2 ступ. \\ 3 ступ. & \centering --- & \centering \arraybackslash --- \\
\hline
\raggedright {'Description': ''} & \centering PereklyuchatelVkh2boolconst1S & \centering Выведено \\ Введено & \centering --- & \centering \arraybackslash --- \\
\hline
\raggedright {'Description': ''} & \centering NezavisimoeRVnasrabatyvanie3Tvremya & \centering 0 ... 100000 & \centering  & \centering \arraybackslash 1 \\
\hline
\raggedright {'Description': ''} & \centering Formirovatelimppoperednemufrontu1Tvremya & \centering 0 ... 100000 & \centering  & \centering \arraybackslash 1 \\
\hline
\raggedright {'Description': ''} & \centering PereklyuchatelVykh3boolconst1S & \centering Откл. \\ Эхо \\ Эхо+ЛОСК & \centering --- & \centering \arraybackslash --- \\
\hline
\raggedright {'Description': ''} & \centering NezavisimoeRVnavozvrat2Tvremya & \centering 0 ... 100000 & \centering  & \centering \arraybackslash 1 \\
\hline
%===m
\end{longtable}
\normalsize

\item
Для работы логики ТУ ДЗ необходимо, чтобы выбранные ступени ДЗ с обоих концов ЛЭП выявили повреждение. Однако в энергосистеме могут быть такие режимы, когда ток повреждения недостаточен для срабатывания защиты, например, короткое замыкание на конце защищаемой линии со слабым питанием. В таком случае не будет пуска ступени ДЗ как на срабатывание, так и на выдачу разрешающего сигнала, что приведёт к отказу логики ТУ ДЗ. Для устранения этого недостатка возможно применение логики отключения слабого конца (ЛОСК) с формированием эхо-сигнала.

\item
Эхо-сигнал формируется с выдержкой времени \textbf{«Тср эхо ТУ ДЗ»} (выдержка времени предназначена для отстройки от разновременного пуска защит по концам линии), если в течение 200 мс отсутствует пуск защит (МФТО, ТНЗНП, ДЗ, АМТЗ).

\item
Длительность эхо-сигнала определяется выдержкой времени \textbf{«Тимп эхо ТУ ДЗ»}. Формирование нового эхо-сигнала блокируется на время, задаваемое уставкой \textbf{«Тбл эхо ТУ ДЗ»}, для исключения зацикливания логики.

\item
Эхо-сигнал действует на выдачу разрешающего сигнала \textit{«Выход ТУ ДЗ»}.

\item
Эхо-сигнал обеспечивает быстрое отключение линии со стороны «сильной» системы (S2), в некоторых случаях это может привести к увеличению тока со стороны «слабой» системы (S1) -- обеспечится каскадное отключение линии (см. рисунок \ref{pict_lo:logic_echo}). Если же после формирования эхо-сигнала не создаются условия каскадного отключения, то предусмотрена возможность отключения слабого конца по критерию минимального напряжения.

\begin{figure}[H]
\centering
\includegraphics[width=0.55\textwidth,height=0.55\textheight,keepaspectratio]{\fbpath/561.1711 Логика связи/img/logic_echo.pdf}
\captionof{figure}{Пояснение к принципу действия эхо-логики}\label{pict_lo:logic_echo}
\end{figure}

\item
Если сформирован эхо-сигнал, то сигнал на отключение формируется при снижении одного из контролируемых фазных напряжений ниже значения уставки \textbf{«Uср ЛОСК»} с контролем включенного положения выключателя. При снижении напряжения во всех фазах действие на отключение разрешено первые 100 мс для ограничения длительности отключающего воздействия. Критерий минимального напряжения блокируется при срабатывании БНН. Функциональная схема логики контроля минимального напряжения приведена на рисунке \ref{pict_lo:logic_LOSK}.

\vspace{2mm}
\begin{figure}[H]
\centering
\includegraphics[width=0.9\textwidth,height=0.9\textheight,keepaspectratio]{\fbpath/561.1711 Логика связи/img/logic_LOSK.pdf}
\captionof{figure}{Функциональная схема логики контроля минимального напряжения ЛОСК}\label{pict_lo:logic_LOSK}
\end{figure}

\item
Режим работы эхо-логики и ЛОСК для ТУ ДЗ определяется положением программного переключателя «Режим эхо ТУ ДЗ», имеющим три положения:
\begin{enumerate}[label=\arabic*)]
    \item   \textit{Откл.} -- эхо-логика и ЛОСК выведены из работы;
    \item   \textit{Эхо} -- в работу введена только эхо-логика;
    \item   \textit{Эхо + ЛОСК} -- в работу введены и эхо-логика, и ЛОСК.
\end{enumerate}

\item
Уставки контроля минимального напряжения ЛОСК представлены в таблице \ref{desc_lo:tbl4}.

\small
\begin{longtable}{|>{\centering\arraybackslash}m{5.3cm}|>{\centering\arraybackslash}m{3.3cm}|>{\centering\arraybackslash}m{4.2cm}|>{\centering\arraybackslash}m{1.8cm}|>{\centering\arraybackslash}m{1cm}|}
\caption{Параметры для настройки контроля минимального напряжения ЛОСК\hfill\vspace{-0.5\baselineskip}}\label{desc_lo:tbl4}\\ 
\hline
\rowcolor{gray!30}
Параметр & Условное обозначение на схеме & Значение/ Диапазон & Единица измерения & Шаг \\ 
\hline
\endfirsthead
\caption*{\hspace{3pt}\emph{Продолжение таблицы \ref{desc_lo:tbl4}\hfill\vspace{-0.5\baselineskip}}} \\ % сделано по ГОСТ 2.105 п.6.8.7
\hline
\rowcolor{gray!30}
Параметр & Условное обозначение на схеме & Значение/ Диапазон & Единица измерения & Шаг \\ 
%\hline
\endhead

%===m>Логика связи|Контроль минимального напряжения ЛОСК
\raggedright {'Description': ''} & \centering Relesktomvozvrataminfloat1xsr & \centering 0,000 ... 100,000 & \centering  & \centering \arraybackslash 0,001 \\
\hline
%===m
\end{longtable}
\normalsize

\end{enumerate}

%%%%%%%%%%%%%%%%%%%%%%%%%%%%%%%%%%%%%%%%%%%%%%%%%%%%%%%%%%%%%%%%%%%%%%  ТЕЛЕУСКОРЕНИЕ ТЗНП %%%%%%%%%%%%%%%%%%%%%%%%%%%%%%%%%%%%%%%%%%%%%%%%%%
\item Телеускорение ступеней ТНЗНП
\begin{enumerate}[label=\arabic{section}.\arabic{subsection}.\arabic{enumi}.\arabic*, labelsep=4pt, leftmargin=0em, itemindent=65pt, parsep=0pt]

\item
Прием сигнала ТУ ТНЗНП осуществляется с контролем срабатывания ИО ТНЗНП выбранной ступени совместно с РНМНП-Р -- выбор ступени ИО ТНЗНП осуществляется программным переключателем 
\textbf{«Ступень ИО ТНЗНП»}.

\item
Для выявления режима реверса мощности используется РНМНП-Б. Выдержка времени \textbf{«Тфикс РНМНП-Б»} обеспечивает фиксацию устойчивого срабатывания РНМНП-Б при КЗ на параллельной линии. Его значение должно быть меньше полного времени отключения повреждения на параллельной линии. Срабатывание РНМНП-Б продлевается на время \textbf{«Тпрод РНМНП-Б»}. По одновременному наличию продленного сигнала срабатывания РНМНП-Б и появлению срабатывания ИО ТНЗНП выбранной ступени совместно с РНМНП-Р фиксируется режим реверса мощности. Минимальное время действия блокировки при реверсе мощности определяется выдержкой времени \textbf{«Тв блок. ТУ ТНЗНП»}. 

\item
Выдержка времени на срабатывание ТУ ТНЗНП регулируется уставкой \textbf{«Тср ТУ ТНЗНП»}.

\item
Разрешающий сигнал \textit{«Выход ТУ ТНЗНП»} на противоположный конец линии выдается с контролем срабатывания ИО ТНЗНП выбранной ступени совместно с РНМНП-Р с контролем реверса мощности.

\item
Действие на отключение и выдача разрешающего сигнала могут быть заблокированы при выявлении броска намагничивающего тока. Для этого предусмотрен программный ключ \textbf{«БНТ ТУ ТНЗНП»}.

\item
Для ТУ ТНЗНП также возможно применение логики отключения слабого конца с формированием эхо-сигнала.

\item
Эхо-сигнал формируется с выдержкой времени \textbf{«Тср эхо ТУ ТНЗНП»}, если в течение 200 мс отсутствует пуск защит (МФТО, ТНЗНП, ДЗ, АМТЗ).

\item
Длительность эхо-сигнала определяется выдержкой времени \textbf{«Тимп эхо ТУ ТНЗНП»}. Фор\-ми\-ро\-ва\-ние нового эхо-сигнала блокируется на время, задаваемое уставкой \textbf{«Тбл эхо ТУ ТНЗНП»}, для исключения зацикливания логики.

\item
Эхо-сигнал действует на выдачу разрешающего сигнала \textit{«Выход ТУ ТНЗНП»}.

\item
Функциональная схема логики ТУ ТНЗНП приведена на рисунке \ref{pict_lo:tu_tnznp}.

\vspace{2mm}
\begin{figure}[H]
\centering
\includegraphics[width=0.9\textwidth,height=0.9\textheight,keepaspectratio]{\fbpath/561.1711 Логика связи/img/logic_TU_TNZNP.pdf}
\captionof{figure}{Функциональная схема ТУ ТНЗНП}\label{pict_lo:tu_tnznp}
\end{figure}

\item
Если после формирования эхо-сигнала не создаются условия каскадного отключения, то предусмотрено отключение слабого конца по критерию минимального напряжения.

\item
Режим работы эхо-логики и ЛОСК для ТУ ТНЗНП определяется положением программного переключателя \textbf{«Режим эхо ТУ ТНЗНП»}, имеющим три положения:

\begin{enumerate}[label=\arabic*)]
    \item   \textit{Откл.} -- эхо-логика и ЛОСК выведены из работы;
    \item   \textit{Эхо} -- в работу введена только эхо-логика;
    \item   \textit{Эхо + ЛОСК} -- в работу введены и эхо-логика, и ЛОСК.
\end{enumerate}

\item
Уставки ТУ ТНЗНП представлены в таблице \ref{desc_lo:tbl5}.

\small
\begin{longtable}{|>{\centering\arraybackslash}m{5.3cm}|>{\centering\arraybackslash}m{3.3cm}|>{\centering\arraybackslash}m{4.2cm}|>{\centering\arraybackslash}m{1.8cm}|>{\centering\arraybackslash}m{1cm}|}
\caption{Параметры для настройки ТУ ТНЗНП\hfill\vspace{-0.5\baselineskip}}\label{desc_lo:tbl5}\\ 
\hline
\rowcolor{gray!30}
Параметр & Условное обозначение на схеме & Значение/ Диапазон & Единица измерения & Шаг \\ 
\hline
\endfirsthead
\caption*{\hspace{3pt}\emph{Продолжение таблицы \ref{desc_lo:tbl5}\hfill\vspace{-0.5\baselineskip}}} \\ % сделано по ГОСТ 2.105 п.6.8.7
\hline
\rowcolor{gray!30}
Параметр & Условное обозначение на схеме & Значение/ Диапазон & Единица измерения & Шаг \\ 
%\hline
\endhead

%===m>Логика связи|ТУ ТНЗНП
\raggedright {'Description': ''} & \centering NezavisimoeRVnasrabatyvanie3Tvremya & \centering 0 ... 100000 & \centering  & \centering \arraybackslash 1 \\
\hline
\raggedright {'Description': ''} & \centering PereklyuchatelVkh7boolconst1S & \centering Откл. \\ 1 ступ. \\ 2 ступ. \\ 3 ступ. \\ 4 ступ. \\ 5 ступ. \\ 6 ступ. & \centering --- & \centering \arraybackslash --- \\
\hline
\raggedright {'Description': ''} & \centering Programmnyjklyuch1s & \centering Введено \\ Выведено & \centering --- & \centering \arraybackslash --- \\
\hline
\raggedright {'Description': ''} & \centering NezavisimoeRVnavozvrat2Tvremya & \centering 0 ... 100000 & \centering  & \centering \arraybackslash 1 \\
\hline
\raggedright {'Description': ''} & \centering NezavisimoeRVnavozvrat1Tvremya & \centering 0 ... 100000 & \centering  & \centering \arraybackslash 1 \\
\hline
\raggedright {'Description': ''} & \centering NezavisimoeRVnasrabatyvanie1Tvremya & \centering 0 ... 100000 & \centering  & \centering \arraybackslash 1 \\
\hline
\raggedright {'Description': ''} & \centering PereklyuchatelVkh2boolconst1S & \centering Выведено \\ Введено & \centering --- & \centering \arraybackslash --- \\
\hline
\raggedright {'Description': ''} & \centering NezavisimoeRVnasrabatyvanie2Tvremya & \centering 0 ... 100000 & \centering  & \centering \arraybackslash 1 \\
\hline
\raggedright {'Description': ''} & \centering Formirovatelimppozadnemufrontu1Tvremya & \centering 0 ... 100000 & \centering  & \centering \arraybackslash 1 \\
\hline
\raggedright {'Description': ''} & \centering PereklyuchatelVykh3boolconst1S & \centering Откл. \\ Эхо \\ Эхо+ЛОСК & \centering --- & \centering \arraybackslash --- \\
\hline
\raggedright {'Description': ''} & \centering NezavisimoeRVnavozvrat4Tvremya & \centering 0 ... 100000 & \centering  & \centering \arraybackslash 1 \\
\hline
%===m
\end{longtable}
\normalsize

\end{enumerate}

\end{enumerate}